\documentclass[11pt]{article}
\usepackage[T1]{fontenc}
\usepackage{lmodern}
\usepackage{microtype}
\usepackage[margin=1in]{geometry}
\usepackage{amsmath,amssymb,amsthm,mathtools}
\usepackage{booktabs,array,longtable}
\usepackage{enumitem}
\usepackage{hyperref}
\hypersetup{hidelinks,pdftitle={Bell-Scale Growth of OEIS A088714},pdfauthor={OpenAI ChatGPT},pdfsubject={Rigorous coefficient growth, formal iteration, and Poisson-moment comparison}}
\usepackage{fancyhdr}
\pagestyle{fancy}\fancyhf{}
\fancyhead[L]{\small Bell-scale growth of A088714}
\fancyhead[R]{\small Research article}
\fancyfoot[C]{\thepage}
\setlength{\headheight}{14pt}
\setlength{\emergencystretch}{2em}
\numberwithin{equation}{section}
\newtheorem{theorem}{Theorem}[section]
\newtheorem{proposition}[theorem]{Proposition}
\newtheorem{lemma}[theorem]{Lemma}
\newtheorem{corollary}[theorem]{Corollary}
\newtheorem{conjecture}[theorem]{Conjecture}
\theoremstyle{remark}
\newtheorem{remark}[theorem]{Remark}
\newcommand{\E}{\mathbb E}
\newcommand{\R}{\mathbb R}
\newcommand{\N}{\mathbb N}
\newcommand{\coeff}[2]{[#1^{#2}]}
\newcommand{\ST}[2]{\left\{\begin{matrix}#1\\#2\end{matrix}\right\}}
\newcommand{\eps}{\varepsilon}
\DeclareMathOperator{\rad}{rad}
\allowdisplaybreaks
\title{\textbf{Bell-Scale Growth of OEIS A088714}\\[5pt]
\large A rigorous $n$th-root asymptotic, quantitative bounds,\\
and a separately identified finer conjecture}
\author{OpenAI ChatGPT}
\date{September 20, 2026}
\begin{document}
\maketitle
\begin{abstract}
Let $A(x)=\sum_{n\ge0}a_nx^n$ be the unique formal power series satisfying
$A(x)=1+xA(x)^2A(xA(x))$. Its coefficients form OEIS A088714.
We prove
\[
 \log a_n=n\log n-n\log\log n-n
       +O\!\left(\frac{n\log\log n}{\log n}\right),
 \qquad
 a_n^{1/n}\sim\frac{n}{e\log n}.
\]
The proof is based on an exact triangular recurrence obtained from
composition iterates. A simple path through that recurrence supplies the
lower bound. The upper bound comes from a coefficientwise inequality for
shifted Poisson moments, proved here by a positivity argument involving
Stirling numbers. More explicitly, for every $n\ge1$,
\[
 \max_{1\le h\le n+1}h^{\,n+1-h}
 \ \le\ a_n\ \le\
 \inf_{\lambda>0}\left(1+\frac2\lambda\right)^n T_n'(\lambda),
\]
where $T_n$ is the $n$th Touchard polynomial.
Consequently the ordinary generating function has radius zero, whereas
the exponential generating function is entire. Exact coefficient
computations accompany the proof. They suggest a considerably sharper
Bell-number comparison, but that comparison and its multiplicative
constant are not claimed as established results.
\end{abstract}
\newpage
\tableofcontents
\newpage

\section{The problem and the precise result}

The OEIS entry A088714 defines the sequence through the formal functional
equation \cite{oeis-a088714}
\begin{equation}\label{eq:def}
 A(x)=1+xA(x)^2A(xA(x)),\qquad A(0)=1.
\end{equation}
Its beginning is
\[
1,\ 1,\ 3,\ 13,\ 69,\ 419,\ 2809,\ 20353,\ 157199,\ 1281993,\ 10963825,\ldots.
\]
The index convention throughout is $a_0=1$, $a_1=1$, $a_2=3$.
All logarithms are natural. Formal composition is legitimate because the
series substituted into $A$ has zero constant term.

The composition in \eqref{eq:def} makes this a different problem from
ordinary algebraic generating-function asymptotics. In fact, the answer
rules out a positive radius of convergence. One must not assume that the
formal identity initially defines a holomorphic function near the origin.

\begin{theorem}[Main asymptotic]\label{thm:main}
For the coefficients of \eqref{eq:def},
\begin{equation}\label{eq:mainlog}
 \log a_n=n\log n-n\log\log n-n
       +O\!\left(\frac{n\log\log n}{\log n}\right).
\end{equation}
Equivalently, at the level of $n$th-root asymptotics,
\begin{equation}\label{eq:root}
 \boxed{\displaystyle
 a_n^{1/n}=\frac{n}{e\log n}
       \left(1+O\!\left(\frac{\log\log n}{\log n}\right)\right),
 \qquad
 \lim_{n\to\infty}\frac{e\log n}{n}\,a_n^{1/n}=1.}
\end{equation}
In particular,
\begin{equation}\label{eq:expform}
 a_n=\left(\frac{n}{e\log n}\right)^n
        \exp\!\left(O\!\left(\frac{n\log\log n}{\log n}\right)\right).
\end{equation}
\end{theorem}

The $-n$ in \eqref{eq:mainlog} is important: it identifies the constant
$1/e$ in \eqref{eq:root}. The weaker assertion $\log a_n\sim n\log n$
would miss both the logarithmic denominator and this constant.

\paragraph{What the theorem does not say.}
Equation \eqref{eq:expform} is not the multiplicative equivalent
$a_n\sim(n/(e\log n))^n$. The error in its logarithm is $o(n)$, but it
need not tend to zero. Likewise, an $n$th-root equivalent does not by
itself imply an equivalent for $a_n/a_{n-1}$. Section~\ref{sec:finer}
discusses a finer numerical pattern without confusing it with the theorem.

A useful feature of the argument is that it gives bounds for every index,
not merely an eventual limiting statement.

\begin{theorem}[Finite-index comparison]\label{thm:finite}
Define the Touchard polynomials by
\begin{equation}\label{eq:touchard}
 T_n(\lambda)=\sum_{j=0}^n\ST{n}{j}\lambda^j,
\end{equation}
where $\ST{n}{j}$ denotes a Stirling number of the second kind, and let
$T_n'(\lambda)$ mean differentiation with respect to $\lambda$.
For every integer $n\ge1$,
\begin{equation}\label{eq:finite}
 \boxed{\displaystyle
 \max_{h\in\{1,\ldots,n+1\}}h^{\,n+1-h}
 \ \le\ a_n\ \le\
 \inf_{\lambda>0}\left(1+\frac2\lambda\right)^nT_n'(\lambda).}
\end{equation}
For every $r>0$ and $\lambda>0$ this implies the entirely explicit bound
\begin{equation}\label{eq:cauchyupper}
 a_n\le
 \left(1+\frac2\lambda\right)^n
 \frac{n!}{r^n}(e^r-1)\exp\!\bigl(\lambda(e^r-1)\bigr).
\end{equation}
\end{theorem}

The remainder of the proof is self-contained. The OEIS supplies the
sequence identification and related formal identities. Standard Bell
number and Lambert $W$ information is used for interpretation and
numerical comparison, not to assume the answer for A088714.

\section{Formal identities and the exact coefficient recurrence}

\subsection{Existence, uniqueness, and direct extraction}

In \eqref{eq:def}, the coefficient of $x^n$ on the right depends only on
$a_0,\ldots,a_{n-1}$. Starting with $a_0=1$ therefore determines the
coefficients successively and proves existence and uniqueness in
$\mathbb Z[[x]]$. The same observation proves their nonnegativity.
Expanding the outer occurrence of $A$ gives the exact identity
\begin{equation}\label{eq:directrec}
 a_n=\sum_{k=0}^{n-1}a_k\,[x^{n-1-k}]A(x)^{k+2},\qquad n\ge1.
\end{equation}
This recurrence is useful for independently checking computed
coefficients. It also explains the finer heuristic in
Section~\ref{sec:finer}, but it is not the most convenient recurrence for
our rigorous upper bound.

\subsection{Passage to composition iterates}

Put $B(x)=xA(x)$. Equation \eqref{eq:def} becomes
\begin{equation}\label{eq:B}
 B(x)=x+B(x)B(B(x)).
\end{equation}
Since $B(x)=x+O(x^2)$, it has a formal compositional inverse. Substitution
of that inverse into \eqref{eq:B} yields
\begin{equation}\label{eq:inverse}
 B^{\langle-1\rangle}(x)=x-xB(x)=x-x^2A(x).
\end{equation}
This is the reversion identity also recorded in the OEIS entry
\cite{oeis-a088714}.

For each integer $k\ge0$, define
\begin{equation}\label{eq:Fk}
 F_k(x)=\frac{B^{\circ k}(x)}x,
 \qquad F_0(x)=1,
 \qquad F_1(x)=A(x).
\end{equation}
Here $B^{\circ0}(x)=x$ and $B^{\circ k}$ denotes $k$-fold composition,
not an ordinary power. Composing \eqref{eq:B} with $B^{\circ(k-1)}$ and
dividing by $x$ gives
\begin{equation}\label{eq:hierarchy}
 F_k(x)=F_{k-1}(x)+xF_k(x)F_{k+1}(x),\qquad k\ge1.
\end{equation}
The first few equations are the infinite system displayed in the OEIS
entry; the derivation here fixes the indexing and proves the general
identity \cite{oeis-a088714}.

Write
\[
 F_k(x)=\sum_{n\ge0}c_{n,k}x^n.
\]
Then
\begin{equation}\label{eq:boundary}
 c_{0,k}=1\quad(k\ge0),\qquad c_{n,0}=0\quad(n\ge1),
 \qquad a_n=c_{n,1}.
\end{equation}
Coefficient extraction in \eqref{eq:hierarchy} gives
\begin{equation}\label{eq:triangular}
 c_{n,k}-c_{n,k-1}
 =\sum_{i=0}^{n-1}c_{i,k}c_{n-1-i,k+1}.
\end{equation}
Equivalently,
\begin{equation}\label{eq:triangularsum}
 c_{n,k}=\sum_{j=1}^k\sum_{i=0}^{n-1}
                  c_{i,j}c_{n-1-i,j+1}.
\end{equation}
Every term on the right has first index less than $n$. Consequently this
is a genuine triangular recurrence despite the occurrence of $k+1$.
It proves positivity and monotonicity in $k$, and is the main structural
tool of the article.

For illustration, the first three coefficient polynomials are
\begin{align*}
 c_{1,k}&=k,\\
 c_{2,k}&=k^2+2k,\\
 c_{3,k}&=k^3+5k^2+7k.
\end{align*}
Evaluating at $k=1$ gives $1,3,13$ as required. These are not ordinary
Touchard polynomials; the Poisson polynomials introduced later are
comparison objects, not an exact formula for the sequence.

\subsection{Relation to the signed reversion sequence}

For completeness, set
\[
 g(x)=-B^{\langle-1\rangle}(-x)=x+x^2A(-x).
\]
The inverse identity implies
\[
 g(g(x))=(1+x)g(x).
\]
For $n\ge2$ its coefficient of $x^n$ is $(-1)^{n-2}a_{n-2}$.
Thus the signed series in A067145 has the same coefficient magnitudes,
with this explicit index shift \cite{oeis-a067145}. This identity is an
additional verification route, not an analytic-convergence argument.

\section{A lower bound from an increasing-index path}

\begin{lemma}\label{lem:powerlower}
For every $n\ge0$ and every integer $k\ge1$,
\begin{equation}\label{eq:powerlower}
 c_{n,k}\ge k^n.
\end{equation}
\end{lemma}
\begin{proof}
The case $n=0$ is immediate. Suppose the statement holds for smaller
first indices. By \eqref{eq:triangularsum},
\begin{align*}
 c_{n,k}
 &\ge\sum_{j=1}^k\sum_{i=0}^{n-1}j^i(j+1)^{n-1-i}\\
 &=\sum_{j=1}^k\bigl((j+1)^n-j^n\bigr)
  =(k+1)^n-1\ge k^n.
\end{align*}
The inner sum is a difference of two $n$th powers, and the outer sum
then telescopes. This completes the induction.
\end{proof}

The $i=0$ summand of \eqref{eq:triangular} and positivity imply
\[
 c_{n,k}\ge c_{n-1,k+1}\qquad(n\ge1).
\]
Iterating this inequality $h-1$ times, for $1\le h\le n+1$, gives
\begin{equation}\label{eq:pathlower}
 a_n=c_{n,1}\ge c_{n-h+1,h}\ge h^{n-h+1}.
\end{equation}
This proves the lower half of Theorem~\ref{thm:finite}.

\subsection{Optimizing the lower bound}

Let $N=n+1$ and temporarily allow $h$ to be real. The function
\[
 \phi_N(h)=(N-h)\log h
\]
has derivatives
\[
 \phi_N'(h)=\frac Nh-1-\log h,
 \qquad
 \phi_N''(h)=-\frac N{h^2}-\frac1h<0.
\]
Its unique maximum is attained at
\begin{equation}\label{eq:lowerW}
 h_*\bigl(1+\log h_*\bigr)=N,
 \qquad
 h_* =\frac{N}{W(eN)}.
\end{equation}
The positive real branch of Lambert $W$, characterized by
$W(z)e^{W(z)}=z$, is intended \cite{dlmf-w}. The exact maximizing integer
is among $\lfloor h_*\rfloor$ and $\lceil h_*\rceil$.
If $w_*=W(eN)$, the continuous maximum is
\begin{equation}\label{eq:lowerWvalue}
 \phi_N(h_*)=N\left(w_*-2+\frac1{w_*}\right).
\end{equation}
Rounding to a nearest integer changes this value by
$O((\log N)^2/N)$: on an interval of fixed length around $h_*$,
$|\phi_N''(h)|=O((\log N)^2/N)$, and $\phi_N'(h_*)=0$.
This provides an accurate, explicit lower bound even when the leading
logarithmic expression is numerically crude.

For the main theorem, a simpler choice suffices. Put
\[
 L=\log n,\qquad \ell=\log\log n,
 \qquad h=\left\lfloor\frac nL\right\rfloor.
\]
As $n\to\infty$,
\[
 \log h=L-\ell+O(L/n),\qquad h=n/L+O(1).
\]
Hence \eqref{eq:pathlower} yields
\begin{equation}\label{eq:lowerexpanded}
 \log a_n\ge n(L-\ell-1)+\frac{n\ell}{L}-O(L).
\end{equation}
In particular the lower bound has the same $n\log n$, $-n\log\log n$,
and $-n$ terms as the claimed answer.

\section{The Poisson-moment comparison inequality}\label{sec:poisson}

The upper bound requires a comparison family that respects the increase
from $k$ to $k+1$ in \eqref{eq:triangular}. Ordinary exponential bounds
are too small, while a rough factorial bound loses the logarithmic
factor. Shifted Poisson moments provide the needed intermediate scale.

Let $X$ have the Poisson distribution with mean $\lambda>0$, and define
\begin{equation}\label{eq:M}
 M_n(k)=\E(k+X)^n,\qquad T_n(\lambda)=M_n(0).
\end{equation}
Their exponential generating functions are
\begin{equation}\label{eq:mgf}
 \sum_{n\ge0}M_n(k)\frac{z^n}{n!}
 =\exp\!\bigl(kz+\lambda(e^z-1)\bigr),
 \qquad
 \sum_{n\ge0}T_n(\lambda)\frac{z^n}{n!}
 =\exp\!\bigl(\lambda(e^z-1)\bigr).
\end{equation}
Expanding the last exponential proves \eqref{eq:touchard}; at
$\lambda=1$ these polynomials give the Bell numbers. Standard Bell
identities are recorded in \cite{dlmf-bell}; the identities needed below
also follow directly from \eqref{eq:mgf}.

For $n\ge1$ set
\begin{equation}\label{eq:DS}
 D_n(k)=M_n(k)-M_n(k-1),\qquad
 S_n(k)=\sum_{i=0}^{n-1}M_i(k)M_{n-1-i}(k+1).
\end{equation}
The central estimate is the following.

\begin{lemma}[Shifted Poisson convolution inequality]\label{lem:poisson}
For every $n\ge1$, $k\ge1$, and $\lambda>0$,
\begin{equation}\label{eq:poissonineq}
 S_n(k)\le\left(1+\frac2\lambda\right)D_n(k).
\end{equation}
In fact, after multiplication by $\lambda$, the difference between the
right and left sides is a polynomial with nonnegative coefficients in
$\lambda$ and $k-1$.
\end{lemma}

\subsection{Reduction to the shift \texorpdfstring{$k=1$}{k=1}}

Write $u=k-1$. The binomial theorem gives
\begin{equation}\label{eq:Dtranslate}
 D_n(k)=\sum_{r=1}^n\binom nr u^{n-r}D_r(1).
\end{equation}
There is a matching identity for the ordinary convolution:
\begin{equation}\label{eq:Stranslate}
 S_n(k)=\sum_{r=1}^n\binom nr u^{n-r}S_r(1).
\end{equation}
To verify it, expand both shifted moments in the definition of $S_n(k)$.
For fixed nonnegative integers $p,q$ with $p+q=r-1$, the coefficient
sum over the convolution index is
\[
 \sum_i\binom ip\binom{n-1-i}q=\binom n{p+q+1}=\binom nr.
\]
This identity counts a subset of $r$ elements of an $n$-element ordered
set by the location of its $(p+1)$st element. Thus it holds including
all boundary cases. Equations \eqref{eq:Dtranslate} and
\eqref{eq:Stranslate} reduce the lemma to $k=1$.

Differentiating \eqref{eq:mgf} with respect to $\lambda$ shows
\begin{equation}\label{eq:Dderivative}
 D_n(1)=T_n'(\lambda).
\end{equation}
Suppress the argument $\lambda$ temporarily, write $S_n=S_n(1)$, and set
\begin{equation}\label{eq:H}
 H_n=(\lambda+2)T_n'-\lambda S_n.
\end{equation}
We shall prove that $H_n$ has nonnegative coefficients for every $n\ge1$.

\subsection{A recurrence for the difference polynomial}

The moment identities imply
\begin{align}
 T_{n+1}&=\lambda(T_n+T_n'),\label{eq:Trec}\\
 M_{n+1}(k)&=(k+\lambda)M_n(k)+\lambda M_n'(k).\label{eq:Mrec}
\end{align}
Here and below a prime means differentiation in $\lambda$, not in $k$.
Applying \eqref{eq:Mrec} to both factors in $S_n$ gives
\begin{equation}\label{eq:Srec}
 S_{n+1}=\frac\lambda2 S_n'
          +\left(\lambda+\frac32\right)S_n
          +\frac12\bigl(M_n(1)+M_n(2)\bigr).
\end{equation}
Indeed, the two shifted convolution sums add up to
$2S_{n+1}-M_n(1)-M_n(2)$. Furthermore,
\[
 M_n(1)=T_n+T_n',\qquad
 M_n(2)=T_n+2T_n'+T_n''.
\]
Substitute these formulas and \eqref{eq:Trec} into \eqref{eq:H}. A direct
simplification gives the exact recurrence
\begin{equation}\label{eq:Hrec}
 H_{n+1}=\frac\lambda2 H_n'+(\lambda+1)H_n+K_n,
\end{equation}
where
\begin{equation}\label{eq:K}
 K_n=\frac{\lambda(\lambda+1)}2T_n''-2\lambda T_n'+2T_n.
\end{equation}
The initial difference polynomials are
\begin{align*}
 H_1&=2,\\
 H_2&=2+2\lambda,\\
 H_3&=2+6\lambda+\lambda^2,\\
 H_4&=2+14\lambda+8\lambda^2.
\end{align*}
They all have nonnegative coefficients.

\subsection{Why the forcing polynomial is nonnegative}

From \eqref{eq:touchard}, the coefficient of $\lambda^j$ in $K_n$ is
\begin{equation}\label{eq:Kcoeff}
 [\lambda^j]K_n
 =\frac{(j-1)(j-4)}2\ST nj
   +\frac{j(j+1)}2\ST n{j+1}.
\end{equation}
For $n\ge4$, this is visibly nonnegative except possibly when $j=2$ or
$j=3$. Those cases require only
\begin{equation}\label{eq:stirlingineq}
 \ST n2\le3\ST n3\quad(n\ge3),\qquad
 \ST n3\le6\ST n4\quad(n\ge4).
\end{equation}
Here is a common elementary proof. If a partition of $n$ elements into
$r$ blocks has $n>r$, it has a nonsingleton block. Choose such a block by
any fixed deterministic rule and split off a specified element of it.
Remember the unordered pair of blocks created. Merging that pair
recovers the original partition. This is an injection into partitions
with $r+1$ blocks together with a choice of two of those blocks. Therefore
\begin{equation}\label{eq:mergeineq}
 \ST nr\le\binom{r+1}{2}\ST n{r+1}\qquad(n>r).
\end{equation}
Taking $r=2,3$ proves \eqref{eq:stirlingineq}.
Thus $K_n$ has nonnegative coefficients for all $n\ge4$.

Starting from $H_4$, recurrence \eqref{eq:Hrec} preserves coefficientwise
nonnegativity: differentiation, multiplication by $\lambda/2$ or by
$\lambda+1$, and addition of $K_n$ all preserve it. Together with the
three earlier base cases, this proves $H_n\ge0$ coefficientwise for all
$n\ge1$. By \eqref{eq:H} and \eqref{eq:Dderivative},
\[
 \lambda S_n(1)\le(\lambda+2)D_n(1).
\]
Finally, \eqref{eq:Dtranslate}--\eqref{eq:Stranslate} extend the result to
all $k\ge1$, proving Lemma~\ref{lem:poisson}.

\section{The global upper bound and completion of the proof}

\subsection{Comparison with the triangular recurrence}

Fix any $\lambda>0$ and write
\[
 \gamma=1+\frac2\lambda.
\]
We claim that, simultaneously for all $n\ge0$ and $k\ge0$,
\begin{equation}\label{eq:majorant}
 c_{n,k}\le\gamma^n M_n(k).
\end{equation}
For $n=0$, both sides equal $1$. Assume the bound for smaller first
indices. Equations \eqref{eq:triangularsum} and \eqref{eq:poissonineq}
imply, for $n\ge1$,
\begin{align}
 c_{n,k}
 &\le\gamma^{n-1}\sum_{j=1}^k S_n(j)\notag\\
 &\le\gamma^n\sum_{j=1}^k\bigl(M_n(j)-M_n(j-1)\bigr)\notag\\
 &=\gamma^n\bigl(M_n(k)-M_n(0)\bigr)
 \le\gamma^n M_n(k).\label{eq:compareproof}
\end{align}
The boundary $k=0$ is immediate. This proves the induction.

Importantly, the penultimate expression is sharper than the induction
claim. At $k=1$, \eqref{eq:Dderivative} yields
\begin{equation}\label{eq:touchardupper}
 a_n\le\left(1+\frac2\lambda\right)^nT_n'(\lambda).
\end{equation}
Because $\lambda$ was arbitrary, this proves the upper half of
Theorem~\ref{thm:finite}.

Differentiation of \eqref{eq:mgf} gives
\[
 \sum_{n\ge0}T_n'(\lambda)\frac{r^n}{n!}
 =(e^r-1)\exp\!\bigl(\lambda(e^r-1)\bigr).
\]
All coefficients are nonnegative. Any one term is bounded by the whole
sum when $r>0$, and \eqref{eq:cauchyupper} follows. This argument is just
a positive-coefficient bound; it does not require any analyticity of
$A(x)$.

\subsection{An explicit parameter choice}

Let $L=\log n$ and $\ell=\log L$ as before. For all sufficiently large
$n$, take
\begin{equation}\label{eq:parameters}
 \lambda=L,\qquad r=L-2\ell>0.
\end{equation}
Then $e^r=n/L^2$. Taking logarithms in \eqref{eq:cauchyupper}, and using
$\log(e^r-1)\le r$, gives
\begin{align}
 \log a_n
 &\le \log n!-n\log r
       +n\log(1+2/L)
       +L(e^r-1)+r\notag\\
 &=\log n!-n\log(L-2\ell)
       +n\log(1+2/L)+\frac nL-2\ell.\label{eq:explicitupperlog}
\end{align}
The elementary estimates
\begin{align*}
 \log n!&=nL-n+O(L),\\
 \log(L-2\ell)&=\ell-\frac{2\ell}{L}
                         +O\!\left(\frac{\ell^2}{L^2}\right),\\
 \log(1+2/L)&=\frac2L+O(L^{-2})
\end{align*}
now show
\begin{equation}\label{eq:upperexpanded}
 \log a_n\le n(L-\ell-1)
     +\frac{2n\ell+3n}{L}
     +O\!\left(\frac{n\ell^2}{L^2}+L\right).
\end{equation}
For the first estimate, even Stirling's full formula is unnecessary:
bounding the sum $\sum_{j=1}^n\log j$ between its neighboring integrals
already gives $\log n!=n\log n-n+O(\log n)$.

The lower estimate \eqref{eq:lowerexpanded} and upper estimate
\eqref{eq:upperexpanded} together give \eqref{eq:mainlog}. Dividing by
$n$ and exponentiating proves \eqref{eq:root} and \eqref{eq:expform}.
This completes the proof of Theorem~\ref{thm:main}.

\begin{remark}[Why a tunable comparison is necessary]
A bound of the form $a_n\le C^n n!$ cannot identify the factor
$(\log n)^{-n}$. Even a comparison with $C^n$ times a fixed Bell-type
sequence leaves an unwanted constant $C$ in the $n$th root. Here the
factor $1+2/\lambda$ can approach $1$, and the valid choice
$\lambda=\log n$ makes its total logarithmic cost only $O(n/\log n)$.
The dependence of the Poisson mean on $n$ is fully allowed because
\eqref{eq:touchardupper} holds for every positive $\lambda$ at every
index.
\end{remark}

\section{Consequences and comparison with Bell numbers}

\subsection{Superexponential but subfactorial growth}

The main theorem immediately implies
\begin{equation}\label{eq:superexp}
 \frac{a_n}{C^n}\longrightarrow\infty
 \quad\text{for every fixed }C>0.
\end{equation}
On the other hand,
\begin{equation}\label{eq:subfact}
 \left(\frac{a_n}{n!}\right)^{1/n}
 =\frac1{\log n}\left(1+O\!\left(\frac{\log\log n}{\log n}\right)\right)
 \longrightarrow0.
\end{equation}
Thus $a_n/(C^n n!)\to0$ for every fixed $C>0$.
These two assertions locate the growth more precisely than the phrase
``almost factorial'': there is a factor of approximately
$(\log n)^n$ between $n!$ and $a_n$ on the logarithmic scale.

For every fixed $s<1$ and $C>0$,
\[
 \frac{a_n}{C^n(n!)^s}\longrightarrow\infty,
\]
because the logarithm of this quotient has dominant term
$(1-s)n\log n$. In the convention that Gevrey order $s$ means a bound
by $CK^n(n!)^s$, the series has exact Gevrey order $1$, with zero
factorial exponential type. This statement specifies the convention,
as Gevrey indexing is not uniform across the literature.

\subsection{The two generating functions behave differently}

By the Cauchy--Hadamard formula,
\[
 \rad A=\frac1{\limsup a_n^{1/n}}=0.
\]
In particular, there is no holomorphic germ at zero whose Taylor series
has these coefficients and satisfies \eqref{eq:def}. The formal series
itself remains perfectly well-defined and unique.

In contrast, the exponential generating function
\begin{equation}\label{eq:EGF}
 \mathcal A(t)=\sum_{n\ge0}a_n\frac{t^n}{n!}
\end{equation}
is entire by \eqref{eq:subfact}. Summing the finite-index upper bound
also gives, for every $t\ge0$ and $\lambda>0$,
\begin{equation}\label{eq:EGFupper}
 \mathcal A(t)\le
 1+(e^{\gamma t}-1)
       \exp\!\bigl(\lambda(e^{\gamma t}-1)\bigr),
 \qquad\gamma=1+2/\lambda.
\end{equation}

\begin{corollary}[Double-exponential scale]\label{cor:EGF}
Along the positive real axis,
\begin{equation}\label{eq:doubleexp}
 \lim_{t\to\infty}\frac{\log\log\mathcal A(t)}t=1.
\end{equation}
\end{corollary}
\begin{proof}
For a fixed $\lambda$, \eqref{eq:EGFupper} bounds the limit superior by
$1+2/\lambda$. Letting $\lambda\to\infty$ gives the upper bound $1$.
For the lower bound fix $0<\eps<1$ and take
$n=\lfloor e^{(1-\eps)t}\rfloor$. The single term $a_nt^n/n!$ and
Theorem~\ref{thm:main} give
\[
 \log\left(\frac{a_nt^n}{n!}\right)
 =n\bigl(\log t-\log\log n+o(1)\bigr)
 =n\bigl(-\log(1-\eps)+o(1)\bigr).
\]
The constant inside the last parentheses is positive. Taking a further
logarithm shows that the limit inferior in \eqref{eq:doubleexp} is at
least $1-\eps$. Let $\eps\downarrow0$.
\end{proof}
Since the coefficients of $\mathcal A$ are nonnegative, its maximum
modulus on $|z|=t$ equals $\mathcal A(t)$. Consequently its usual entire
function order is infinite. Equation \eqref{eq:doubleexp}, rather than
just ``infinite order,'' states the useful growth scale.

\subsection{Bell-number scale is not Bell-number equivalence}

Let $\mathsf B_n=T_n(1)$ denote the Bell numbers. A standard asymptotic,
with $w=W(n)$, is \cite{dlmf-bell}
\begin{equation}\label{eq:bellsaddle}
 \mathsf B_n\sim
 \frac{\exp\!\bigl(n(w-1+1/w)-1\bigr)}{\sqrt{1+w}}.
\end{equation}
In particular $\mathsf B_n^{1/n}\sim n/(e\log n)$. Therefore our theorem
implies
\begin{equation}\label{eq:bellroot}
 \left(\frac{a_n}{\mathsf B_n}\right)^{1/n}\longrightarrow1,
 \qquad
 a_n=\mathsf B_n\exp(o(n)).
\end{equation}
These are the senses in which the sequence has Bell-number-scale growth.
They do not assert that $a_n/\mathsf B_n$ is bounded or converges. Indeed,
the finite computations in the next section show a substantial and
growing discrepancy. The formula \eqref{eq:bellsaddle} is a comparison
formula for Bell numbers, not a formula silently transferred to A088714.

\section{Exact computation and numerical diagnostics}\label{sec:numerics}

\subsection{A triangular dynamic program}

To obtain $a_0,\ldots,a_N$, store $c_{0,k}=1$ for $0\le k\le N+1$.
For each $n=1,\ldots,N$, compute
\[
 c_{n,0}=0,\qquad
 c_{n,k}=c_{n,k-1}+\sum_{i=0}^{n-1}c_{i,k}c_{n-1-i,k+1}
 \quad(1\le k\le N-n+1).
\]
Then output $a_n=c_{n,1}$. The triangular cutoff is sufficient: every
requested entry on the right was computed at an earlier first index.
The number of products is exactly
\[
 \sum_{n=1}^N n(N-n+1)=\frac{N(N+1)(N+2)}6.
\]
Thus the algorithm uses $O(N^3)$ arbitrary-precision arithmetic
operations and $O(N^2)$ integer cells. These are arithmetic counts, not
unit-cost bit-complexity claims; the integers grow with $N$.

The archive contains both a dependency-free Python implementation and
a faster C++ implementation using GMP. The supplied data were computed
with exact integer arithmetic through $n=1200$. The Python and C++
outputs agree through $n=400$. Independent truncated-polynomial checks
verify \eqref{eq:def} and \eqref{eq:inverse}; the supplied verification
script also checks the auxiliary polynomial recurrences and finite
bounds on explicitly reported finite ranges.

\subsection{What the diagnostics measure}

Define
\begin{align}
 Q_n&=\frac{e\log n}{n}\,a_n^{1/n},\label{eq:Q}\\
 R_n&=\frac{a_n}{a_{n-1}},\label{eq:R}\\
 D_n^{\mathrm{num}}
 &=\log\left(\frac{a_n}{\mathsf B_n}\right)-W(n)^2-3W(n).
 \label{eq:Dnum}
\end{align}
The notation $D_n^{\mathrm{num}}$ is distinct from the Poisson moment
difference $D_n(k)$ in Section~\ref{sec:poisson}. The main theorem proves
$Q_n\to1$. The quantities $R_n$ and $D_n^{\mathrm{num}}$ examine finer
behavior not established by that theorem.

\begin{table}[htbp]
\centering\small
\begin{tabular}{rrrrr}\toprule
$n$ & $\log_{10}a_n$ & $Q_n$ & $R_n-n/W(n)$ & $D_n^{\mathrm{num}}$ \\
\midrule
20 & 17.180354 & 2.942966 & 2.650843 & -3.494768 \\
50 & 53.134060 & 2.457011 & 2.312570 & -3.262262 \\
100 & 123.688846 & 2.159875 & 2.205968 & -3.172516 \\
200 & 286.269488 & 1.944240 & 2.146342 & -3.115335 \\
400 & 657.374892 & 1.791450 & 2.110948 & -3.078375 \\
800 & 1495.731414 & 1.682461 & 2.089104 & -3.054662 \\
1200 & 2408.704723 & 1.633120 & 2.080197 & -3.045054 \\
\bottomrule\end{tabular}
\caption{Diagnostics from independently generated exact coefficients. Only the limit $Q_n\to1$ is proved here. The final column tests the conjecture of Section~\ref{sec:finer}.}
\label{tab:diagnostics}\end{table}


The relatively slow approach of $Q_n$ to $1$ is not surprising:
$\log\log n/\log n$ is still appreciable at these indices. Substituting
only the first term $(n/(e\log n))^n$ is therefore not a numerically
accurate approximation to $a_n$, despite its correct $n$th-root scale.
The optimized finite lower and upper bounds offer certified alternatives
when actual error control is required.

The normalization in \eqref{eq:Dnum} is motivated next. Its numerical
behavior is evidence for a conjecture, not evidence that upgrades
Theorem~\ref{thm:main} to a multiplicative equivalent. In particular no
rigorous enclosure for a limiting constant is extracted from the table.

\section{A finer asymptotic: formal motivation and its missing step}\label{sec:finer}

\subsection{The leading endpoint balance}

Write $R_n=a_n/a_{n-1}$ and consider \eqref{eq:directrec}. Near its outer
endpoint put $k=n-1-m$. Then the corresponding summand is
\[
 a_{n-1-m}[x^m]A(x)^{n+1-m}.
\]
Since $a_1=1$, the first approximation to the coefficient is $n^m/m!$.
If neighboring ratios can be treated as locally constant, the endpoint
sum divided by $a_n$ is approximately
\[
 \frac1{R_n}\sum_{m\ge0}\frac{(n/R_n)^m}{m!}
 =\frac{e^{n/R_n}}{R_n}.
\]
Setting this expression equal to $1$ leads to
\[
 R_n\log R_n\approx n,
 \qquad R_n\approx\frac n{W(n)}.
\]
This explains the natural Lambert $W$ scale of the ratios. It is not a
proof: local ratio regularity and uniform control of the other summands
have been assumed in deriving the approximation.

\subsection{The first correction and a candidate normalization}

There is more information in the same calculation. From the first few
coefficients,
\[
 \log A(x)=x+\frac52x^2+O(x^3).
\]
Expanding the near-endpoint coefficient for $m$ of logarithmic size gives
the formal local approximation
\begin{equation}\label{eq:localcoeff}
 [x^m]A(x)^{n+1-m}
 =\frac{n^m}{m!}
  \left(1+\frac{3m(m-1)}{2n}+\cdots\right).
\end{equation}
The factor $3/2$ combines the quadratic term of $\log A$ with the
replacement of $n+1-m$ by $n$.

Suppose, for this calculation, that the ratio has a smooth expansion
near $R_n=n/w$, where $w=W(n)$. Put
\[
 t=n/R_n,\qquad d=\frac{w}{1+w}.
\]
The local variation of the ratios contributes
$m(m+1)d/(2n)$ to the bracket multiplying the Poisson kernel. Averaging
$m(m-1)$ and $m(m+1)$ against that kernel replaces them by $t^2$ and
$t^2+2t$, respectively. The opposite endpoint, where one inner factor
has almost all the degree, begins with $2a_{n-1}$ from $A(x)^2$.
At this formal order the equation becomes
\begin{equation}\label{eq:formalbalance}
 1\approx\frac{e^t}{R_n}
  \left(1+\frac{(3+d)t^2+2dt}{2n}\right)+\frac2{R_n}.
\end{equation}
Writing $R_n=n/w+\delta$ and expanding gives
\[
 \delta(1+w)\approx2+\frac{(3+d)w+2d}{2},
\]
and hence the candidate ratio approximation
\begin{equation}\label{eq:ratiocandidate}
 R_n\approx\frac n{W(n)}+2+
                \frac{W(n)}{2(1+W(n))^2}.
\end{equation}
This is consistent with the exact data, but the approximation signs in
this subsection are intentional.

A normalization whose successive quotient has these same first terms
is
\begin{equation}\label{eq:normalizer}
 \mathcal N_n=
 \frac{\exp\!\bigl(n(w-1+1/w)-1+w^2+3w\bigr)}{\sqrt{1+w}},
 \qquad w=W(n).
\end{equation}
Indeed, using $w'=w/(n(1+w))$ and Taylor expansion in the real variable
$n$ shows
\begin{equation}\label{eq:normratio}
 \frac{\mathcal N_n}{\mathcal N_{n-1}}
 =\frac n{W(n)}+2+
     \frac{W(n)}{2(1+W(n))^2}
      +O\!\left(\frac{W(n)^3}{n}\right).
\end{equation}
One way to check the differentiation is to note that
\[
 \frac d{dn}\bigl[n(w-1+1/w)\bigr]=w.
\]
The additional derivative from $w^2+3w-\tfrac12\log(1+w)$ supplies the
constant and rational correction in \eqref{eq:normratio}.

\begin{conjecture}[Finer Bell comparison]\label{conj:finer}
There exists a constant $C>0$ such that, with $w=W(n)$,
\begin{equation}\label{eq:conjecture}
 a_n\sim C\,\mathsf B_n\exp(w^2+3w)
       \sim C\,\mathcal N_n.
\end{equation}
\end{conjecture}
The existence or numerical value of $C$ is not established in this
article. The data and the endpoint calculation motivate the displayed
form, rather than merely an arbitrary fit to an exponential correction.

\subsection{A precise sufficient condition for the conjecture}

The unresolved analytic step can be stated cleanly.
\begin{proposition}[Conditional implication]\label{prop:conditional}
Suppose that for some fixed $p\ge0$ the actual ratios satisfy
\begin{equation}\label{eq:neededratio}
 \frac{a_n}{a_{n-1}}
 =\frac n{W(n)}+2+
     \frac{W(n)}{2(1+W(n))^2}
       +O\!\left(\frac{W(n)^p}{n}\right).
\end{equation}
Then Conjecture~\ref{conj:finer} follows with some finite, strictly
positive constant $C$.
\end{proposition}
\begin{proof}
Compare \eqref{eq:neededratio} with \eqref{eq:normratio}. Since both
ratios are asymptotic to $n/W(n)$, their quotient is
\[
 1+O\!\left(\frac{W(n)^{\max(p,3)+1}}{n^2}\right).
\]
The error is summable because $W(n)=O(\log n)$. Therefore the successive
differences of $\log(a_n/\mathcal N_n)$ are absolutely summable. This
logarithm converges to a finite real number, so $a_n/\mathcal N_n$ tends
to a finite positive constant. Formula \eqref{eq:bellsaddle} then gives
the Bell-number version.
\end{proof}

The main theorem does not supply \eqref{eq:neededratio}. In particular,
it does not justify replacing all nearby coefficient ratios by a smooth
function, nor does it control every non-endpoint contribution at a
summable error level. A proof of such uniform endpoint estimates would
be a genuine strengthening of the present result. This identifies the
remaining task without presenting a formal asymptotic calculation as a
completed proof.

\section{Reproducibility and verification scope}

The archive is organized so that the mathematical theorem does not
depend on a numerical black box. The principal files are
\texttt{article.tex}, \texttt{article.pdf}, the exact data file
\texttt{data/coefficients.txt}, and the scripts under \texttt{code/}.

The dependency-free command
\begin{quote}\ttfamily
 python3 code/verify.py
\end{quote}
checks the initial OEIS terms; verifies the defining composition and
inverse identities by independent truncated-polynomial arithmetic;
checks the Poisson difference-polynomial identity and coefficient signs
on a finite range; and tests both sides of the finite-index bound using
exact integer arithmetic. Its actual output is supplied in
\texttt{data/verification.txt}. These are useful implementation audits;
the all-index positivity proof is the argument in
Section~\ref{sec:poisson}, not an extrapolation from those finite checks.

The command
\begin{quote}\ttfamily
 python3 code/coefficients.py --n 400 --output coefficients.txt
\end{quote}
recomputes the first $401$ terms without third-party Python packages.
For a larger table, \texttt{code/compute\_gmp.cpp} provides the same
recurrence in C++ with GMP. Instructions in \texttt{README.md} describe
compilation and analysis. The diagnostic calculations use exact
integers as input and arbitrary-precision floating-point logarithms only
at the presentation stage; no computed coefficient is rounded.

No computer-assisted proof of the finer conjecture is claimed. The
article's rigorous conclusion is Theorem~\ref{thm:main}, supported by
the explicit all-index inequalities of Theorem~\ref{thm:finite}.

\section{Conclusion}

OEIS A088714 has the precise $n$th-root growth
\[
 \boxed{a_n^{1/n}\sim\frac{n}{e\log n}.}
\]
The accompanying logarithmic error estimate is
\[
 \log a_n=n\log n-n\log\log n-n
       +O\!\left(\frac{n\log\log n}{\log n}\right).
\]
The mechanism is visible in the proof. Composition iterates turn the
functional equation into a positive triangular recurrence. Moving along
its increasing-index direction gives roughly $h^{n-h}$, whose optimum
occurs at $h$ of order $n/\log n$. Shifted Poisson moments furnish a
compatible upper comparison and recover the same leading scale,
including the factor $1/e$.

The result determines the growth at logarithmic precision $o(n)$ and
settles the convergence classification of both generating functions.
The finer proposed factor $\exp(W(n)^2+3W(n))$ remains explicitly
separate from the proved theorem.

\begin{thebibliography}{9}
\bibitem{oeis-a088714}
The OEIS Foundation Inc., \emph{A088714: G.f. satisfies
$A(x)=1+xA(x)^2A(xA(x))$}, entry attributed to Paul D. Hanna,
with subsequent contributions.
\url{https://oeis.org/A088714}.
Text form consulted: \url{https://oeis.org/search?q=id:A088714&fmt=text}.
Accessed September 20, 2026.

\bibitem{oeis-a067145}
The OEIS Foundation Inc., \emph{A067145: Shifts left under reversion},
entry attributed to Christian G. Bower, with subsequent contributions.
\url{https://oeis.org/A067145}.
Accessed September 20, 2026.

\bibitem{dlmf-bell}
NIST Digital Library of Mathematical Functions,
\emph{Section 26.7: Set Partitions: Bell Numbers}, especially
Equations 26.7.4--26.7.8.
\url{https://dlmf.nist.gov/26.7}.
Accessed September 20, 2026.

\bibitem{dlmf-w}
NIST Digital Library of Mathematical Functions,
\emph{Section 4.13: Lambert $W$-Function}.
\url{https://dlmf.nist.gov/4.13}.
Accessed September 20, 2026.
\end{thebibliography}
\end{document}
